\SourceLeaf{160}but one should make it a rule to conceal nothing, so that others can calculate differently if they wish. The numerical values of the results will be equally serviceable however one proceeds. But one risks overstating their reliability if one is too quick to exclude observations. Decisions of this kind seem to me more akin to \textls{acting in everyday life}, where mathematical rigour and certainty are seldom or never available and one must be content to proceed according to one's best considered judgement.

\begin{center}
. . . . . . . . . . . . . . . . .
\end{center}



\Needspace{6\baselineskip}
\begin{center}
\textsc{Gauss} to \textsc{Gerling}. Göttingen, 2 April 1840.
\end{center}

. . . . . . . . . . As to the scientific content of your letter, I shall merely say a few words in reply to the difficulty you raise, unnecessarily, on account of the heterogeneity of angular and linear measurements.

Suppose $q, q', q'', q'''$ etc. are quantities found by observation (whether homogeneous or heterogeneous makes no difference) and that $dq, dq', dq''\ldots$ denote the differences between observation and calculation according to some system of elements. The principle of the method of least squares \textls{by no means} requires that
\[
dq^2+dq'^2+dq''^2+\text{etc.}
\]
be a minimum. It requires that
$\left(\dfrac{dq}{m}\right)^2+\left(\dfrac{dq'}{m'}\right)^2+\left(\dfrac{dq''}{m''}\right)^2+$etc.
be a minimum, where $m, m', m''\ldots$ denote the respective mean errors to be expected in the individual data. Cases in which $m=m'=m''=m'''$ etc. can or must be assumed occur so overwhelmingly often in practice that one is led to accept the first statement and to forget that only the second is generally valid.

To make use of this principle one must have at least an approximate knowledge, where exact knowledge is unavailable, if not of the quantities
